\scanpage{497}
\mspage{551}

% 497-01
A priori, $\alpha$ and $\beta$ depend on the choice of $u_0$ from which $u$ comes
not only on $u$ --- but
% 497-02
\[
(16)\qquad
\left\{\begin{aligned}
\alpha\rho(\alpha)^{-1}&=[u,\rho]\\
\beta\sigma(\beta)^{-1}&=[u,\sigma]
\end{aligned}\right.
\]
% 497-03
do not depend on it. If we put
% 497-04
\[
(17)\qquad
T_{\rho}=\uline{\mathrm{Centr}}_M(\rho)
\qquad
T_{\sigma}=\uline{\mathrm{Centr}}_M(\sigma)
\]
% 497-05
we thus see that, modulo $T_{\rho}$ resp. $T_{\sigma}$,
$\alpha,\beta$ depend only on $u$.

% 497-06, 497-07, 497-08, 497-09
\uline{Hypothesis}
$\alpha\inserted{{}=\alpha(u_0)},\beta\inserted{{}=\beta(u_0)}$
depend only on $u=\varphi(u_0)$
(for $u_0\in\deleted{\unsure{M}}\unsure{G}[0]$).

% 497-10
We shall \unsure{express} this hypothesis in a slightly different way,
by introducing
% 497-11
\[
(18)\qquad J=\operatorname{Ker}(G[0]\longrightarrow M[0]),
\]
% 497-12, 497-13
we thus want that for $\deleted{\unsure{\alpha}}$
% 497-14, 497-15
\[
u'_0=\deleted{u_0}\delta_0u_0\quad
\text{with }u_0\in G[0],\ \delta_0\in J,\quad\text{we have}
\]
% 497-16
\[
\alpha(u'_0)=\alpha(u_0),\qquad \beta(u'_0)=\beta(u_0)
\]
% 497-17, 497-18
\unsure{Write}
\[
\begin{gathered}
\alpha_0=\uline{\alpha}(u_0),\quad
\beta_0=\uline{\beta}(u_0),\quad
\alpha'_0=\uline{\alpha}(u'_0),\quad
\beta'_0=\uline{\beta}(u'_0),\quad
\overline{\alpha}_0=\uline{\alpha}(\delta_0),\quad
\overline{\beta}_0=\uline{\beta}(\unsure{\delta_1|\delta_0}),
\end{gathered}
\]
\uline{in} $G$, so that we shall have, by the cocycle formulas of $G$
% 497-19
\[
\alpha'_0=\delta_0(\alpha_0)\overline{\alpha}_0,
\qquad
\beta'_0=\delta_0(\beta_0)\overline{\beta}_0,
\]
% 497-20
whence, applying $\varphi$ and noting that
% 497-21
\[
\begin{aligned}
\varphi(\delta_0(\alpha_0))
&=\underbrace{\varphi(\delta_0)}_{1}(\varphi(\alpha_0))
=\varphi(\alpha_0)\overset{\mathrm{dfn}}=\alpha
&&\text{and }\overline{\mathrm m}\\
\varphi(\delta_0(\beta_0))&=\varphi(\beta_0)\overset{\mathrm{dfn}}=\beta
\end{aligned}
\]
% 497-22, 497-23
\[
(19)\qquad
\alpha'=\deleted{\unsure{\varphi}}\alpha\overline{\alpha},\qquad
\beta'=\beta\overline{\beta}\qquad\text{in }\unsure{\mathcal{C}}
\]
% 497-24, 497-25
where $\alpha,\beta,\alpha',\beta',\overline{\alpha},\overline{\beta}$ are the images under
$\varphi$ of the quantities $\alpha_0$ etc in $\unsure{\mathcal{C}}_{0,3}^{+\wedge}$.
\deleted{\unsure{with} $\alpha=\alpha(u_0)$}
Thus the hypothesis means this
% 497-26
\[
(20)\qquad
\forall\,\delta_0\in J\overset{\mathrm{dfn}}=
\operatorname{Ker}(G[0]\longrightarrow M[0]),\quad
\text{we have}\quad
\unsure{\alpha^{\Pi}}(\delta_0),\unsure{\beta^{\Pi}}(\delta_0)
\in\unsure{\operatorname{Ker}}\varphi
\]
% 497-27, 497-28
\deleted{\unsure{when}}
When this hypothesis is satisfied, we thus have maps
\mspage{496}

\scanpage{498}

% 498-01
\[
(21)\qquad
\diagram{\begin{tikzcd}[column sep=large]
M[0]\arrow[r,shift left=0.65ex,"\uline{\alpha}"]
\arrow[r,shift right=0.65ex,"\uline{\beta}"']&\unsure{\mathcal{C}}
\end{tikzcd}}
\]
% 498-02, 498-03
\deleted{satisfying}
\unsure{factoring} (14), and thus satisfying
% 498-04
\[
(22)\qquad
\left\{\begin{aligned}
u(\rho)&=\operatorname{int}(\alpha)\rho\\
\unsure{u}(\sigma)&=\operatorname{int}(\beta)\sigma
\end{aligned}\right.
\qquad
\text{for }\unsure{u}\in M[0],\quad
\alpha=\uline{\alpha}(u),\quad\beta=\uline{\beta}(u),
\]
% 498-05, 498-06, 498-07
which \deleted{\unsure{generalize|generalizing}}
\inserted{\unsure{paraphrase}} (2').

% 498-08
Moreover, let us put
% 498-09
\[
(23)\qquad
\unsure{\zeta|\xi}=\alpha\rho(\alpha)^{-1},\qquad
\eta=\beta\sigma(\beta)^{-1}
\]
% 498-10
the relations (22) can be written thus
% 498-11
\[
(24)\qquad u(\rho)=\unsure{\zeta|\xi}\rho,\qquad u(\sigma)=\eta\sigma
\]
% 498-12
and we shall have
% 498-13
\[
(25)\qquad \unsure{\zeta|\xi}=[u,\rho],\qquad \eta=[u,\sigma]
\]
% 498-14
and moreover
% 498-15
\[
(26)\qquad
\unsure{\zeta|\xi}=\eta\unsure{\varepsilon_1|\ell_1}^{\nu}\qquad
\text{where }\nu=\frac{\mu-1}{2}\in\widehat{\mathbb Z}
\quad\text{($\mu$ the multiplier of $u$).}
\]
% 498-16
which comes from the single relation in $\unsure{\mathcal{C}}_{0,3}^{+}$.

% 498-17
(We assume to fix ideas that the homomorphism
% 498-18, 498-19
\[
\deleted{\unsure{(27)}}\qquad
\diagram{\begin{tikzcd}[column sep=large]
G\arrow[r,"\chi^{G}"]&\widehat{\mathbb Z}^{*}
\end{tikzcd}}
\qquad\text{(``multiplier'')}
\]
% 498-20
factors through
% 498-21
\[
(27)\qquad
\diagram{\begin{tikzcd}[column sep=large]
G\arrow[r]&M\arrow[r,"\chi^{M}"]&\widehat{\mathbb Z}^{*},
\end{tikzcd}}
\]
% 498-22
so that we have a notion of multiplier ($\in\widehat{\mathbb Z}^{*}$)
for the $u\in M$.)

% 498-23
Consider the natural homomorphism
% 498-24
\[
(28)\qquad
\diagram{\begin{tikzcd}[column sep=large,row sep=small]
M[0]\arrow[r]&\operatorname{Aut}(\unsure{\mathcal{C}}):\\
u\arrow[r,mapsto]&\operatorname{int}(u)|\unsure{\mathcal{C}}
\end{tikzcd}}
\]
% 498-25, 498-26
we see from (22) \inserted{\unsure{even|or even} (24)}
that it is known when we know
\mspage{497}

\scanpage{499}
\mspage{552}

% 499-01, 499-02
\margin{for \deleted{\illegible{1}} $\rho,\sigma$ generate
\unsure{topologically} $\unsure{\mathcal{C}}$}
% 499-03
the functions $\uline{\alpha},\uline{\beta}:M[0]\longrightarrow\unsure{\mathcal{C}}$,
and even when we know only
$\unsure{\zeta|\xi},\eta:M[0]\longrightarrow\unsure{\mathcal{C}}$.
If moreover we put
% 499-04
\[
(29)\qquad
M'[0]=\operatorname{Im}(M[0]\longrightarrow\operatorname{Aut}(\unsure{\mathcal{C}}))
\]
% 499-05
these latter functions $\unsure{\zeta|\xi},\eta$ are already defined on $M'[0]$ --- whereas it
is not clear that the same holds for $\alpha,\beta$, when
$M[0]$ \unsure{does not act|does not operate} faithfully on $\unsure{\mathcal{C}}$.
We thus have a map \inserted{$(\unsure{\zeta|\xi},\eta)$} from a certain subgroup
$M'[0]\subset\operatorname{Norm}_{\operatorname{Aut}(\unsure{\mathcal{C}})}(L_0)$
to the set $\unsure{\mathcal{R}}$ of solutions of the equation in $\unsure{\zeta|\xi},\eta$
(26) which we \unsure{write|shall write} simply
% 499-06
\[
(30)\qquad \eta^{-1}\unsure{\zeta|\xi}\in L_1
\]
% 499-07, 499-08
a relation which implies (26) (with the precise form of the exponent $\nu$)
if we assume \inserted{(\unsure{for example})} $\unsure{\mathcal{C}|M}$ big enough that
% 499-09
\inserted{$\diagram{\begin{tikzcd}[column sep=large]
\unsure{\mathcal{C}}_{0,3}^{+\wedge}\arrow[r,"\unsure{2}"']&
\unsure{\mathrm{SL}|\mathrm{PSL}}(2,\widehat{\mathbb Z})
\end{tikzcd}}$}
factors through $\unsure{\mathcal{C}}$ --- which we shall assume, to make life easier for ourselves.

% 499-10, 499-11
\deleted{We shall make on $\unsure{\mathcal{C}}$, equipped with $\rho,\sigma,\varepsilon_1$
the following hypothesis}
\inserted{More precisely, let}
% 499-12
\deleted{(31) \unsure{H)|II)} On}
% 499-13, 499-14, 499-15, 499-16
\[
(32)\qquad
\left\{\begin{gathered}
\deleted{\unsure{\Sigma|\mathcal{R}}}\ \unsure{=}
\Sigma_{\unsure{\mathcal{C}},\rho,\sigma}
\subset\unsure{\mathcal{C}}\times\unsure{\mathcal{C}}\\
\Vert\\[-1ex]
\left\{(\unsure{\zeta|\xi},\eta)\ \middle|
\begin{array}{ll}
\text{a)}&\unsure{\zeta|\xi}\rho\ \text{conjugate to }\rho\quad
\deleted{\text{in }\operatorname{Aut}(\unsure{\mathcal{C}})}
\inserted{\text{in }\unsure{\mathcal{C}}}\\[1ex]
\text{b)}&\eta\sigma\quad\rule{2.7cm}{0.4pt}\quad\sigma\quad\rule{2cm}{0.4pt}\\[1ex]
\text{c)}&\eta^{-1}\unsure{\zeta|\xi}\in L_1\quad\text{i.e.}\quad
\exists\nu'\in\widehat{\mathbb Z}\ \text{such that }\unsure{\zeta|\xi}=\eta\varepsilon_1^{\nu'}
\end{array}\right\}
\end{gathered}\right.
\]
% 499-17, 499-18
\margin{\[
\begin{gathered}
\underbrace{[\alpha,\rho]=\alpha\rho(\alpha)^{-1}}\\
\left\{\begin{array}{l}
\text{\unsure{i.e.} }\unsure{\zeta|\xi}\ \text{\unsure{of the form} }\unsure{[\alpha,\rho]}\\
\text{\unsure{resp.} }\eta\ \text{\unsure{of the form} }\unsure{[\beta,\sigma]}
\unsure{{}=\beta\sigma(\beta)^{-1}}
\end{array}\right.
\end{gathered}
\]}

% 499-19, 499-20, 499-21, 499-22, 499-23
\deleted{For \unsure{every} $u\in\operatorname{Aut}(\unsure{\mathcal{C}})$
\uline{normalizing} $L_0$, \deleted{and \unsure{fixing}}
\deleted{\unsure{must}} and fixing the conjugacy class of
$\rho$ and that of $\sigma$, i.e. \unsure{such|such} that
\[
\deleted{\unsure{\zeta|\xi}=[u,\rho]}\qquad(\unsure{\zeta|\xi},\eta)=([u,\rho],[u,\sigma])
\]
\unsure{satisfy|satisfies} the first two conditions of (33), \unsure{and}}
\mspage{498}

\scanpage{500}

% 500-01
and let $u\in\operatorname{Aut}(\unsure{\mathcal{C}})$ be an automorphism of
$\unsure{\mathcal{C}}$, and let us put
% 500-02
\[
(34)\qquad
\left\{\begin{aligned}
\unsure{\zeta|\xi}&=[u,\rho]=u\rho u^{-1}\rho^{-1}=u(\rho)\rho^{-1}
&&\text{i.e. }u(\rho)=\unsure{\zeta|\xi}\rho\\
\eta&=[u,\sigma]=u(\sigma)\sigma^{-1}
&&\text{i.e. }u(\sigma)=\eta\sigma
\end{aligned}\right.
\]
% 500-03
\margin{cf. (24),(25)}
% 500-04, 500-05, 500-06
then to say that $\unsure{\zeta|\xi}\rho$ \inserted{$(=u(\rho))$} is conjugate to $\rho$ means
that $u$ fixes the conjugacy class of $\rho$, to say that $\eta\sigma$
\inserted{$(=u(\sigma))$} is conjugate \unsure{of|to} $\sigma$ means that $u$ fixes
the conjugacy class of $\sigma$.
Finally the condition $\eta^{-1}\unsure{\zeta|\xi}\in L_1$ can also be written, since
% 500-07, 500-08
\[
\begin{aligned}
\eta^{-1}\unsure{\zeta|\xi}
&=(\sigma u(\sigma^{-1}))(u(\rho)\rho^{-1})
=\sigma u(\sigma^{-1}\rho)\rho^{-1}\\
&=\sigma\left(u(\underbrace{\sigma^{-1}\rho}_{\varepsilon_0})
\deleted{\unsure{\sigma}}\underbrace{\rho^{-1}\sigma}_{\varepsilon_0^{-1}}\right)
\end{aligned}
\]
% 500-09
\margin{N.B. $\sigma^{-1}=\sigma$}
% 500-10, 500-11
\deleted{\unsure{whence}}
\[
{}=\sigma(u(\varepsilon_0)\varepsilon_0^{-1})
\]
% 500-12
in the form
% 500-13
\[
\sigma(u(\varepsilon_0)\varepsilon_0^{\unsure{-1}})\in L_1
\qquad\text{i.e.}\qquad
u(\varepsilon_0)\varepsilon_0^{-1}\in L_0
\]
% 500-14, 500-15, 500-16
(since $L_0=\sigma^{-1}(L_1)$) or again
(since $\varepsilon_0$ \deleted{generates}\inserted{$\in L_0$})
% 500-17
\[
u(\varepsilon_0)\in L_0,\quad
\text{i.e. (since $\varepsilon_0$ generates $L_0$)}
\]
% 500-18
\margin{In this argument, we have used only the relation
$\sigma^{-1}\rho=\varepsilon_0$ and the fact that $\varepsilon_0$ generates $L_0$,
\unsure{and not} that $\sigma^2=1$ or $\rho^3=1$ for example.}
% 500-19
\[
u(L_0)=L_0
\]
% 500-20
i.e. that $u$ \uline{normalizes} $L_0$.

% 500-21
Thus let
% 500-22
\[
(35)\qquad
\left\{\begin{gathered}
\widetilde{B}'_0\subset\operatorname{Aut}(\unsure{\mathcal{C}})\\
\Vert\ \mathrm{dfn}\\[-1ex]
\left\{u\in\operatorname{Aut}(\unsure{\mathcal{C}})\ \middle|
\begin{array}{l}
u\ \text{fixes conjugacy \unsure{class} of $\rho$ and that of $\sigma$}\\
u\ \text{normalizes }L_0
\end{array}\right\},
\end{gathered}\right.
\]
% 500-23
so that we have a map (injective because $\rho,\sigma$ generate
$\unsure{\mathcal{C}}$)
% 500-24
\[
(36)\qquad
\diagram{\begin{tikzcd}[column sep=large,row sep=small]
\widetilde{B}'_0\arrow[r,hook]&\Sigma\\
u\arrow[r,mapsto]&([u,\rho]=u(\rho)\rho^{-1},\ [u,\sigma]=u(\sigma)\sigma^{-1}).
\end{tikzcd}}
\]
\mspage{499}

\scanpage{501}
\mspage{553}

% 501-01--07: heading, double marginal rule, hypothesis and its alterations.
\[
(37)\quad
\begin{array}{||l}
\text{\uline{Fundamental hypothesis on} }\unsure{\mathcal{C}},\ 
\text{\uline{equipped with} }\rho,\sigma,\ 
\text{\deleted{and}\inserted{\uline{whence}} }\varepsilon_0=\sigma^{-1}\rho,\\[1ex]
\varepsilon_1=\rho\sigma^{-1}=\sigma(\varepsilon_0)=\rho(\varepsilon_0):\\[1ex]
\text{a) }\unsure{\mathcal{C}}\text{ is generated topologically by }\rho,\sigma\\[1ex]
\text{b) The map (36) \deleted{\unsure{is|by}} (injective by a)) is}\\
\text{\phantom{b) }bijective.}
\end{array}
\]

% 501-08--10: canceled first version and its replacement.
\deleted{This thus means that $\forall(\unsure{\zeta|\xi},\eta)\in\Sigma$ \unsure{is|u}}

Condition b) thus means that for
\deleted{of the form $u$,}
$(\alpha,\beta)\in\unsure{\mathcal{C}}\times\unsure{\mathcal{C}}$,
satisfying a relation
\[
\uline{\alpha\rho(\alpha)^{-1}}
=\uline{\beta\sigma(\beta)^{-1}}\varepsilon_1^{\nu'},
\qquad \unsure{\nu''|\nu'}\text{ suitable},
\]
there exists $u\in\operatorname{Aut}(\unsure{\mathcal{C}})$ such that we have
% 501-12--14: two identities, including the canceled right-hand attempt.
\[
\begin{array}{lll}
[u,\rho]=\alpha\rho(\alpha)^{-1}&
\text{i.e. }u(\rho)=\operatorname{int}(\alpha)(\rho)&
\text{i.e. }\inserted{\unsure{u(\rho)|u\rho)}=\unsure{\zeta|\xi}\rho}
\deleted{\unsure{\operatorname{int}(\unsure{\zeta|\xi})\rho}}\\[1ex]
[u,\sigma]=\beta\sigma(\beta)^{-1}&
\text{i.e. }u(\sigma)=\operatorname{int}(\beta)\sigma&
\text{i.e. }u(\sigma)=\eta\sigma
\end{array}
\]

% 501-15: prose continuation (tilde is not visible on this B).
\margin{$\unsure{\bullet}$}
(we shall thus necessarily have $u\in B'_0\ \unsure{\}}$, and
$u$ will be unique if \unsure{we have} a)).

% 501-16--18: final paragraph and overwritten subscript.
Thus, on the set \unsure{of the|of} $\unsure{R|\mathcal{R}}_{\rho,\sigma}$ of the $(\unsure{\zeta|\xi},\eta)$ satisfying
a) b) c), we shall have a group structure, \unsure{+|and} \unsure{deduced|described} by transport of structure from that of
$\widetilde{B}'_{\deleted{\unsure{\rho}}0}$

\unsure{etc.}

\mspage{500}